% Source pin (SHA-256): PifoStatement.lean=fe9475ebfa2c462d27cb0a99781926074e1281396f2c9e8d32d6dfca788d814e.
\section{The model}
\label{sec:model}

This section presents the mathematical content of the frozen statement file.
In particular, ranks and arrival stamps are natural numbers, topologies are
finite ordered trees, a push carries one global stamp through the entire path,
and a failed pop is observable and atomic.

\subsection{Entries and PIFOs}

Let \(X\) be a set of values.  A queue entry is a triple
\[
  e=(\val(e),\rank(e),\arr(e))\in X\times\Nat\times\Nat.
\]
The value is either a packet type or a child index, depending on the node at
which the entry resides.  Lower ranks have higher priority.  The strict
selection relation is lexicographic:
\begin{equation}
 e\mathrel{\prec}e'
 \quad\Longleftrightarrow\quad
 \rank(e)<\rank(e')
 \ \lor\ 
 \bigl(\rank(e)=\rank(e')\land\arr(e)<\arr(e')\bigr).
 \label{eq:entry-order}
\end{equation}
A PIFO queue is represented as a list in insertion order.  Popping removes the
least entry under \autoref{eq:entry-order} and preserves the relative list order
of the rest.  Every push receives a fresh global arrival stamp and contributes
at most one entry to a particular node.  Consequently arrival stamps within a
reachable queue are distinct, so its minimum is unique.  The arrival component
implements FIFO tie breaking among entries of equal rank.

\subsection{Topologies, states, and paths}

A topology is generated by
\[
  \tau ::= \mathsf{leaf}
      \mid \mathsf{node}[\tau_0,\ldots,\tau_{m-1}],
\]
where the child list is finite, ordered, and possibly empty.  A state over
packet values \(A\) has a queue at every node.  A leaf queue contains entries
whose values lie in \(A\); an internal queue contains entries whose values are
natural-number child indices.  The empty state has the given topology and an
empty queue at every node.

An annotated root-to-leaf path is
\[
  p ::= \mathsf{leaf}(r)
      \mid \mathsf{node}(c,r,p),
  \qquad c,r\in\Nat.
\]
It is well formed for a topology when the constructors match at every level,
each \(c\) is an in-range child index, and the remaining path fits that child.
Pushing packet \(a\) with global arrival stamp \(j\) along a valid path appends
one entry at each visited node:
\begin{align*}
 \mathsf{node}(c,r,p)&:\quad (c,r,j)
     &&\text{is appended to the internal queue, then the push follows child }c,\\
 \mathsf{leaf}(r)&:\quad (a,r,j)
     &&\text{is appended to the leaf queue.}
\end{align*}
Thus every entry created by one source push has the same arrival stamp.  The
formal definition also specifies ill-formed pushes, but the main theorem
quantifies only over valid schedulers, so those branches are unreachable.

To pop a tree, pop its root queue.  At a leaf, the selected value is emitted.
At an internal node, the selected value is interpreted as a child index and the
pop continues recursively in that child.  The operation is \emph{atomic}: if
the root queue is empty, an index is invalid, or any recursive pop fails, the
tree-level result is \(\None\), with no updated state exposed to the caller.
For valid reachable states, the latter two cases do not arise; the balance
argument appears in \lemref{lem:balance}.

\begin{figure}[t]
\centering
\begin{tikzpicture}[
  q/.style={draw,rounded corners,minimum width=2.1cm,minimum height=.72cm,
            align=center,fill=white},
  every edge/.style={draw,-{Stealth[length=2mm]}},
  lab/.style={font=\small,inner sep=1pt},
  node distance=1.35cm and 1.75cm]
  \node[q] (root) {root PIFO\\child indices};
  \node[q,below left=of root] (left) {leaf PIFO\\packets};
  \node[q,below right=of root] (mid) {internal PIFO\\child indices};
  \node[q,below left=of mid] (mleft) {leaf PIFO\\packets};
  \node[q,below right=of mid] (mright) {leaf PIFO\\packets};
  \draw[->] (root) -- node[lab,left] {$c=0$} (left);
  \draw[->,very thick,pifored] (root) -- node[lab,right] {$c=1,\ r_0$} (mid);
  \draw[->,very thick,pifored] (mid) -- node[lab,left] {$c=0,\ r_1$} (mleft);
  \draw[->] (mid) -- node[lab,right] {$c=1$} (mright);
  \node[lab,pifored,below=.14cm of mleft] {leaf rank $r_2$};
  \node[align=left,anchor=west,font=\small] at ($(root.east)+(1.25,0)$)
    {a push appends along one path;\\a pop chooses anew at each level};
\end{tikzpicture}
\caption{A PIFO-tree state.  Internal entries name children; leaf entries name
packets.  The highlighted annotated path installs ranks \(r_0,r_1,r_2\) with
one shared arrival stamp.  A later pop follows whichever child entry is then
minimal; it need not follow a single source push's path.}
\label{fig:model-tree}
\end{figure}

\subsection{Stateless schedulers and observations}

Fix \(k\in\Nat\) and the finite packet alphabet \(A=\Fin(k)\).  A scheduler is
a topology together with an assignment
\[
  S=(\tau,\mathsf{assign}),
  \qquad \mathsf{assign}:A\to\mathsf{Path}.
\]
It is \emph{valid} when \(\mathsf{assign}(a)\) fits \(\tau\) for every packet
type \(a\).  It is \emph{stateless} because every occurrence of type \(a\)
uses exactly the same annotated path: the same child choices and the same rank
at every level.  The two schedulers compared below may have different
topologies.

An operation is either \(\push(a)\) or \(\popop\).  A run starts from the empty
tree and counter zero.  The \(j\)-th push increments the counter to \(j\) and
uses \(j\) as the arrival stamp.  It emits no observation.  Each pop contributes
exactly one element of \(\mathsf{Option}(A)\): \(\Some(a)\) on success and
\(\None\) on failure.  A failed pop leaves both the tree and arrival counter
unchanged.  Definedness is therefore observable.

The observation records a packet's type, not its source-push identity.  In the
stateless model these observations contain the same ordering information.
Every two occurrences of one type carry equal ranks in every queue they share,
so FIFO tie breaking releases them in their source arrival order.  Hence the
\(i\)-th emitted occurrence of type \(a\) is the \(i\)-th pushed occurrence of
that type.  Relative to a fixed input operation word, the type trace therefore
determines the source-occurrence identity trace.

For a word \(w=a_1\cdots a_n\in A^*\), define its flush operation word by
\begin{equation}
  \flushOps(w)
  =\push(a_1)\cdots\push(a_n)\,\popop^n.
  \label{eq:flush-ops}
\end{equation}
Exactly \(n\) pops suffice: fewer reveal a prefix of the full deterministic
drain, while additional pops append only \(\None\).  Write \(\run(S,o)\) for
the list of pop observations produced by operation word \(o\).

\begin{definition}[The two equivalences]
For valid schedulers \(S,T\) on \(A\),
\begin{align*}
  S\flusheq T
  &\quad\Longleftrightarrow\quad
    \forall w\in A^*.\ 
    \run(S,\flushOps(w))=\run(T,\flushOps(w)),\\
  S\intereq T
  &\quad\Longleftrightarrow\quad
    \forall o\in\mathsf{Op}(A)^*.\ 
    \run(S,o)=\run(T,o).
\end{align*}
\end{definition}

The frozen proposition is the following statement, uniformly over the finite
alphabet size.

\begin{theorem}[Flushes suffice for stateless PIFO trees]
\label{thm:main}
For every \(k\in\Nat\) and all valid stateless schedulers \(S,T\) over
\(\Fin(k)\),
\[
  S\flusheq T \quad\Longrightarrow\quad S\intereq T.
\]
Consequently flush and interleaved equivalence coincide on valid stateless
PIFO trees.
\end{theorem}

The implication from right to left in the final sentence follows immediately
because every \(\flushOps(w)\) is an operation word.  The nontrivial implication
is proved in \autoref{sec:main-proof}.

\subsection{Correspondence with the frozen Lean statement}

\autoref{tab:lean-map} maps the paper notation to the statement's definitions.
The source file is statement-only: its final item is a definition of a
proposition, not a proof of that proposition.

\begin{table}[H]
\centering
\small
\begin{tabularx}{\textwidth}{@{}>{\raggedright\arraybackslash}p{.34\textwidth}
  >{\ttfamily\raggedright\arraybackslash}p{.33\textwidth}
  >{\raggedright\arraybackslash}X@{}}
\toprule
Paper object & \normalfont Lean name & Source lines\\
\midrule
Entry and its three fields & Entry; .val, .rank, .arr & 53--57\\
PIFO contents & Queue & 59--60\\
Strict entry comparison & better & 64--66\\
PIFO pop & qpop & 68--77\\
Tree shape & Topology & 81--86\\
Runtime tree state & Tree & 88--94\\
Empty tree and forest & emptyTree; emptyForest & 96--104\\
Annotated path & Path & 106--113\\
Path well-formedness & pathOk; pathOkAt & 115--126\\
Tree push and list helper & treePush; treePushAt & 128--143\\
Atomic tree pop and list helper & treePop; treePopAt & 145--182\\
Stateless scheduler & Scheduler & 186--193\\
Topology and fixed assignment & Scheduler.topo; Scheduler.assign & 191--192\\
Validity & Scheduler.Valid & 194--197\\
Push/pop operations & Op & 216--220\\
Execution from a state & runFrom & 222--235\\
Execution from the empty state & run & 237--240\\
Flush experiment & flushOps & 244--250\\
Flush equivalence & flushEquiv & 252--254\\
Interleaved equivalence & interEquiv & 256--261\\
Main proposition & statelessFlushImpliesInterleaved & 263--269\\
\bottomrule
\end{tabularx}
\caption{Mathematical-to-Lean correspondence for the frozen formal model.}
\label{tab:lean-map}
\end{table}
